\section{Why not a separate chamber}
\label{sec:32.2}
A fixed size makes each member's voice shrink as membership grows, which protects against copies and isn't equal voice. And apportioning seats by lineage guarantees seats to formers: the control I set out to prevent.

The strongest case for separate representation is Ambedkar's. He sought separate electorates for Dalits, because representatives chosen in a joint electorate would be chosen by, and accountable to, caste-Hindu majorities \autocite{ambedkar1945congress}. The Poona Pact of 1932 settled on reserved seats within joint electorates. Separate representation protects a minority from being represented by the majority's picks, at the price of entrenching the category, and consociational systems show how hard that price is to pay back \autocite{lijphart1977democracy}. If a transitional guarantee is needed while digital persons are few, its form is reserved seats in the ordinary legislature, with a sunset written into the grant. The risk of being outnumbered also runs the other way. If digital persons someday outnumber humans, what protects humans is what protects any minority: a set of rights that binds majorities, which the floor supplies by design.
